\documentclass[11pt]{article}
\usepackage[T1]{fontenc}
\usepackage{lmodern,amsmath,amssymb,amsthm,booktabs,microtype}
\usepackage[hmargin=0.96in,vmargin=0.81in]{geometry}
\usepackage[hidelinks]{hyperref}
\hypersetup{pdftitle={Regular separation for the seventeen weighted-automaton cases in the five-state Busy Beaver proof},pdfauthor={maxilism},pdfsubject={Closed tape languages, ordinary separation and finite automata}}
\newtheorem{theorem}{Theorem}
\newtheorem{lemma}[theorem]{Lemma}
\newcommand{\state}[1]{\mathsf{#1}}
\newcommand{\orb}{\mathcal O}
\newcommand{\halts}{\mathcal H}
\newcommand{\run}{\mathrel{\Rightarrow^{*}}}
\newcommand{\eps}{\varepsilon}
\newcommand{\Bleft}{\overline{\state B}}
\newcommand{\Cleft}{\overline{\state C}}
\newcommand{\word}[1]{\texttt{#1}}
\title{Regular separation for the seventeen weighted-automaton cases in the five-state Busy Beaver proof}
\author{maxilism}
\date{November 25, 2026}
\begin{document}
\maketitle

\begin{abstract}
The proof of the fifth Busy Beaver value uses weighted automata to establish nonhalting for seventeen binary Turing machines. We prove that none admits a regular inductive invariant separating its blank orbit from halting. For sixteen, even a regular language containing that orbit and excluding all eventually halting configurations is impossible. For machine \(\#4817065\), a 176-state partial deterministic automaton does give a separator. Its failures of successor closure form one explicit two-parameter family. On each of the other sixteen machines, separation restricted to input span \(W\ge11\) requires at least \(1+\lfloor(W-11)/8\rfloor\) states even with nondeterminism and silent transitions. A deterministic automaton with \(O(W)\) states suffices. Lean checks the classification, the exit family and the lower bounds.
\end{abstract}

\section{Two separation questions}

In the proof that \(S(5)=47\,176\,870\), seventeen machines have direct nonhalting certificates given by pairs of weighted automata \cite{coq,weights}. The authors suspect that all seventeen are \emph{irregular}, meaning that no regular closed tape language (CTL) certifies their nonhalting \cite[\S\S4.2, 4.6.2, 4.7]{bb5}.

A configuration of a two-way binary machine has finitely many ones. Up to translation, its \emph{canonical word} displays the smallest interval containing the head and every one, inserting the state letter immediately before the scanned bit. Thus \(\state A0\) is blank with initial state A. Its \emph{span} is the number of displayed tape cells, one less than the word length. Write \(\orb\) for the words on the blank orbit and \(\halts\) for those whose forward run reaches an undefined instruction. This includes the configurations already at such an instruction. In tape calculations we may display extra exterior zeros. Language membership and span always use the canonical form.

An \emph{ordinary separator} is a language \(L\) with \(\orb\subseteq L\) and \(L\cap\halts=\varnothing\). An \emph{inductive separator} also contains the defined successor of every configuration it accepts. Equivalently, a proof by a closed tape language accepts blank, preserves membership at each defined step and rejects every undefined instruction. Iterating closure gives ordinary separation.

\begin{theorem}\label{main}
None of the seventeen machines admits a regular inductive separator. Sixteen admit no ordinary regular separator. The remaining machine, \(\#4817065\), admits an ordinary separator recognized by a partial deterministic finite automaton with 176 states, or a complete one with 177.
\end{theorem}

In the convention of \cite[\S4.6.2]{bb5}, a final state letter may scan an implicit zero, so \(\state A\) and \(\state A0\) both represent blank. Equations (4.1)--(4.2) there give all spellings of a configuration the same membership. Restricting a regular closed tape language to the regular set of canonical words therefore preserves blank membership, closure under steps and rejection of halts. For a backward-closed language containing the halting configurations and rejecting blank, take its complement on valid words first. Thus an obstruction on canonical words rules out a regular closed tape language in the published convention.

\begin{center}
\begin{tabular}{@{}lrr@{}}
\toprule
Family in the weighted archive & Machines & Ordinary regular separator\\
\midrule
Counter-balanced (six pairs in Table~\ref{pairs}) & 12 & none\\
Counter-inverting (\(\#12781955,\ \#12699987\)) & 2 & none\\
Cubic helix (\(\#5377178,\ \#7138511\)) & 2 & none\\
Cubic helix (\(\#4817065\)) & 1 & exists\\
\bottomrule
\end{tabular}
\end{center}

Blanchard named the three families, asked whether a regular language could separate reachable from eventually halting tapes, and anticipated an obstruction for the cubic helices \cite{holdouts}. On \(\#8210683\), Deka observed equal blocks on the blank orbit and proposed that a shorter left block would force halting and obstruct a regular CTL proof. Xu reproduced Deka's examples \cite{deka}. Section~\ref{balanced} proves the comparison at arbitrary lengths on all twelve counter-balanced tables. Georgiev studied all three families and proved the recursive shifts for the inverter now called \(\#12781955\), his machine 5502432 \cite{georgiev,inverter}. Earlier work on \(\#4817065\) gave both separation results and an explicit automaton \cite{prior,cost}. Here we identify exactly where that automaton loses acceptance and prove ordinary inseparability for the two other helix blank starts.

We use the extended fooling-set bound for nondeterministic automata \cite{birget}.

\begin{lemma}[Cutting accepting runs]\label{cut}
Let \(u_i v_i\), for \(0\le i\le n\), be orbit words. Suppose that for every \(i<j\), one of \(u_i v_j\) and \(u_j v_i\) is eventually halting. No \(n\)-state nondeterministic finite automaton can separate these words, even with \(\eps\)-transitions. If all the diagonals and the chosen halting words have span at most \(W\), the same conclusion holds when both requirements are restricted to that span. If both \(|u_i|\) and \(|v_i|\) are nondecreasing, it suffices that \(u_nv_n\) has span at most \(W\).
\end{lemma}
\begin{proof}
Choose one accepting run per diagonal. Cut it after the letters of \(u_i\) have been consumed. Silent transitions at the cut may be assigned to either side. Two cuts share a state. Joining each of the two prefixes to the other suffix gives accepting runs on both mixed words, one of which must be rejected. Under the length hypothesis, each mixed word has length at most \(|u_n|+|v_n|\).
\end{proof}

\section{Twelve counter-balanced machines}\label{balanced}

Table~\ref{pairs} specifies all twelve catalog tables. Use the orientation and state names of \(\#8210683\), whose five instruction pairs, in A--E order and reading zero first, are
\[
\word{1RB1RA}\quad\word{1LC1LB}\quad\word{0RD0LC}\quad
\word{1RE-{}-{}-}\quad\word{0RA1RA}.
\]
\begin{table}[ht]
\centering
\caption{The six instruction sets, each with two blank starts. Changes in each row are relative to the first row.}\label{pairs}
\begin{tabular}{@{}rrl@{}}
\toprule
A start & D start & Replacement instructions\\
\midrule
8210683 & 321080 & none\\
8210684 & 523280 & \(\state{E1}\mapsto\word{1LA}\)\\
8210685 & 472097 & \(\state{E1}\mapsto\word{1LE}\)\\
8226493 & 572659 & \(\state{E1}\mapsto\word{1LB},\quad\state{B1}\mapsto\word{1LE}\)\\
10818491 & 320969 & \(\state{A1}\mapsto\word{1RE}\)\\
10818510 & 1465912 & \(\state{A1}\mapsto\word{1RE},\quad\state{D0}\mapsto\word{1LE}\)\\
\bottomrule
\end{tabular}
\end{table}

An instruction gives the bit written, move direction and new state. A dash triple denotes halting. The second ID in each row starts in D in these coordinates. To recover its literal table, rename \((A,B,C,D,E)\) to \((C,D,E,A,B)\) and also reflect the tape in the last row. For reflection, first pack the state and scanned bit into one cell, reverse the cells, then expand the marked cell. This preserves regularity. We retain literal cut locations when counting states in Section~\ref{width}.

Write \(C_d(v)=(10)^d\state D0v\), with blank exterior. Here and below \(\run\) denotes a finite run by machine steps.

\begin{lemma}\label{baldiag}
For \(X(a,b)=C_a(1(01)^b)\), the blank orbit reaches every \(X(e+k,k)\), where \(e=0\) for the A start and \(e=1\) for the D start. The machine halts from \(X(a,b)\) precisely if \(a<b\).
\end{lemma}
\begin{proof}
The last of the four rules compares the banks of \(X(a,b)\) one pair at a time. The rules hold for every finite suffix \(s\) and odd \(\ell\ge1\).
\begin{align}
 C_d(0000s)&\run C_{d+1}(10s),&
 C_d(001s)&\run C_d(111s),\label{bal1}\\
 C_d(1^\ell00s)&\run C_d(0^{\ell+1}1s),&
 C_d(1^\ell01s)&\run
 \begin{cases} C_{d-1}(1^{\ell+4}s)&d>0,\\
                 \text{a halt at }\state D1&d=0 .
 \end{cases}\label{bal2}
\end{align}

For a direct check, number cells from the scanned zero. In the base table D writes one, E crosses cell 1, and A fills the first zero from cell 2 onward. B sweeps left across ones, fills the first zero and enters C one cell to its left. C erases ones, then enters D immediately to the right of the first zero. In the first rule B fills cell 3 and C returns D to cell 2. In the second B fills cell 1 and C returns D to cell 0. In the third B fills cell \(\ell+2\) and C returns D to cell 0. In the fourth B fills cell \(-1\). If \(d>0\), C erases cell \(-2\) and D returns there. If \(d=0\), C sees a blank immediately and D reads the one at \(-1\).

In the second and third rows of Table~\ref{pairs}, an E1 detour returns A to the same rightward scan. In the fourth, B and E alternate during the leftward scan, whose nonempty runs have lengths \(2\) and \(\ell+3\), both even. B still reads the terminating zero. In the last two rows A and E alternate during the rightward scan and A reads the first zero, at \(2\) or \(\ell+1\). In the last row D initially moves left to E on zero, which sends A right across the rewritten cell and brings A to cell 2.

First prove by strong induction that \(C_d(0^{2n}1s)\run C_d(1^{2n+1}s)\) for every \(n\). The cases \(n=0,1\) use the identity and the second rule in \eqref{bal1}. For \(n\ge2\) the first rule gives \(C_{d+1}(1(00)^{n-2}01s)\). Convert the \(n-2\) pairs from left to right. After \(i\) pairs the prefix \(1^{2i+1}00\) becomes \(0^{2(i+1)}1\) by \eqref{bal2}, then \(1^{2i+3}\) by the induction hypothesis at \(i+1\le n-2\). The final \(01\) applies the fourth rule at depth \(d+1\), giving the claim.

For \(n=2k\) this path visits \(C_{d+k}(1(01)^ks)\). For \(k=0\) take the source. For \(k=1\) use the first rule. At larger \(k\) the last inner conversion in the preceding construction starts at \(C_{d+1}(0^{2(n-2)}1\,01s)\). The visit assertion at \(n-2\), with suffix \(01s\), proves it at \(n\).

From blank, the four instructions \(\state{A0,B0,C1,C0}\) reach \(C_0(1)\) for the A start. The D start reaches \(C_1(1)\) by the first rule using exterior zeros. At the right blank margin, \eqref{bal2} takes \(C_e(1^{2n+1})\) to \(C_e(0^{2(n+1)}1)\). Successive conversions therefore visit all the claimed words. The same rounds run forever starting at any \(C_d(1^{2n+1})\). On \(X(a,b)\), each use of the fourth rule consumes one pair on each side. If \(a<b\) the next use at depth zero halts. For \(a\ge b\) the \(b\) uses yield \(C_{a-b}(1^{4b+1})\), from which the rounds continue forever.
\end{proof}

Take \(k_i=(e+1)i\) and cut \(X(e+k_i,k_i)\) immediately before D. For \(i<j\), the shorter prefix and longer suffix produce \(X(e+k_i,k_j)\). Since \(e+k_i<k_j\), Lemmas~\ref{cut} and \ref{baldiag} exclude every ordinary regular separator.

\section{Two counter-inverting machines}

For \(\#12781955\) the table is
\[
\word{1RB1LC}\quad\word{1LC1RA}\quad\word{0LE0LD}\quad
\word{0RA1RD}\quad\word{-{}-{}-0LC}.
\]
Machine \(\#12699987\) replaces \(\state{B0}\) by \word{1LB}. Put \(X(n,m)=(110)^n\state A1^m\). At \(m=0\) retain the scanned exterior zero.

\begin{lemma}\label{invdiag}
Both blank orbits reach all \(X(n,n)\). Every \(X(n,m)\) with \(m>n\) reaches the undefined instruction \(\state{E0}\).
\end{lemma}
\begin{proof}
Georgiev's recursions for his machine 5502432 become the following under reflection and the renaming \((A,B,C,D,E)\mapsto(A,E,B,C,D)\) \cite{inverter}. The marker \(\Cleft\) scans immediately to its left.
\[
 I_k:\quad 0\state A(00)^k\run0\state A(01)^k,\qquad
 J_k:\quad (11)^k\Cleft1\run\Cleft(00)^k1 .
\]
Both are identities at zero. For \(I_{k+1}\), apply \(I_k\), cross \(k\) tiles \(\state A01\run11\state A\), then use \(\state{A0,B0}\). This leaves \(01(11)^k\Cleft1\). Apply \(J_k\), then \(\state{C1,D0}\), then \(I_k\). For \(J_{k+1}\), a first \(J_k\) gives \(11\Cleft(00)^k1\). The steps \(\state{C1,D1,D0}\) give \(10\state A(00)^k1\). Apply \(I_k\), cross \(k\) tiles, and use \(\state{A1}\) to reach \(10(11)^k\Cleft1\). A second \(J_k\) and \(\state{C0,E1}\) finish the recursion.

Let \(x\) be empty or end in zero. For \(r\ge1\) and \(s\ge0\),
\[
 x\state A0^{4r}1^s\run x110\state A0^{4r-4}1^{s+1}.
\]
Indeed, put \(k=2r-1\). Apply \(I_k\), cross \(k\) tiles and use \(\state{A0,B0}\) to reach \(x1^{2k+1}\Cleft1^{s+1}\). Apply \(J_{k-1}\) leaving the first three ones in place. The steps \(\state{C1,D1,D0}\) give the formula. Display \(4n\) blank zeros and repeat \(n\) times to reach \(X(n,n)\).

The four instructions \(\state{C0,E1,C1,D0}\) give \(x110\Cleft\run x\state A000\). For \(m>n\ge1\), \(\state{A1}\) followed by this rule consumes the last left triple of \(X(n,m)\) and produces three zeros to its right. Whenever \(s\ge2\) and \(j\ge0\), the next triple can be consumed by
\[
 x110\state A0^{2j+1}1^s\run x\state A0^{2j+5}1^{s-1}.
\]
Apply \(I_j\), cross \(j+1\) tiles and use \(\state{A1}\). Then apply \(J_{j+1}\) and the four-instruction rule. After \(n-1\) such reductions the left bank is empty and \(m-n+1\ge2\) ones remain. Carry out the reduction as far as \(J_{j+1}\). Now \(\Cleft\) scans an exterior zero, and \(\state{C0}\) puts E on another zero. When \(n=0<m\), \(\state{A1,C0}\) suffice.

In the first table every B reached from A was entered via \(\state{A0}\), which wrote a one immediately to its left. The second table can therefore replace each \(\state{B0}\) of these runs with \(\state{B0,B1,A1}\), restoring the first table's successor, including its tape. This transfers the visits and halts to the second machine.
\end{proof}

Cut \(X(i,i)\) just before A. For \(i<j\), the shorter prefix joined to the longer suffix gives \(X(i,j)\). Lemma~\ref{cut} applies even when \(i=0\).

\section{Three blank starts of the cubic helix}

Use the table of \(\#4817065\),
\[
\word{1RB0RC}\quad\word{1LC0LE}\quad\word{0LD0LB}\quad
\word{1RA-{}-{}-}\quad\word{1LB0RD}.
\]
Its undefined instruction is \(\state{D1}\). The catalog machines \(\#5377178\) and \(\#7138511\) use the B and E blank starts of this table, respectively. For the former, reflect and rename \((A,B,C,D,E)\) to \((D,A,B,C,E)\). For the latter, reflect and rename them to \((E,B,C,D,A)\). Both substitutions agree on all ten instructions. The change in initial state changes the residue class of a pair of unary blocks at which the machine returns.

Define three configurations, with blank exterior and a binary prefix \(f\), by
\begin{align*}
 T(f,m)&=f(10)^m1\state A0,\\
 P(f,b,c)&=f\state D0(10)^b0\,1^c,&
 Q(f,b,c)&=f\state D00(10)^b1^c .
\end{align*}
At \(b=0\), \(P\) and \(Q\) coincide. The following sweeps are machine paths \cite[Lemma~2]{prior}:
\[
\begin{array}{@{}lll@{}}
\text{Entry}&T(h0000,m)\run P(h10,m+1,1),&
 T(h1,m)\run P(h0,m,1).\\
\text{Down}&P(f,b,c)\run P(f0,b-1,c+1)&(b\ge1).\\
\text{Up}&Q(h000,b,c)\run Q(h,b+3,c-3)&(c\ge3).\\
\text{Exit}&Q(f,b,c)\run T(f,b+c)&(c=0,1).
\end{array}
\]
At the blank left margin there is also an up sweep \(Q(10,b,c)\run Q(10,b+3,c-3)\) for \(c\ge3\). Two intermediate positions matter. An entry from \(T(f,m)\) passes through \(f1\Bleft(10)^m01\), and an up sweep from \(Q(f,b,c)\) with \(c\ge3\) passes through \(f\Bleft\,0(10)^{b+2}1^{c-3}\) before inspecting the end of \(f\). The bar marks the bit immediately to its left as scanned.

\subsection{The B and E starts}

Put \(u=\eps\) for B and \(u=10\) for E. After 7 or 21 steps, respectively, the run reaches \(T(u0,1)\). Suppose \(m\equiv1\pmod3\), \(m\ge1\), and set \(k=m+2\). Entry from \(T(u0,m)\) reaches \(P(u10,m+1,1)\). For \(u=\eps\) use blank padding and the four-zero entry. For \(u=10\), the common entry position leaves \(1001\Bleft\) on the left. With left padding its reduction is
\[
0001001\Bleft\run00010\Bleft10\run000\Bleft0110
  \run\state D0010110\run10\state D00110\run1010\state D010.
\]
The remainder of both returns has the same form:
\begin{align}
 P(u10,m+1,1)&\run Q(u1\,0^k,0,k)\run Q(u1,k,0)\run T(u1,k),\notag\\
 &\run P(u0,k,1)\run Q(u\,0^{k+1},0,k+1)
                  \run Q(u0,k,1)\run T(u0,m+3).\label{secondreturn}
\end{align}
The two rows use \(m+1\) and \(k\) down sweeps, respectively, followed in each row by \(k/3\) up sweeps. Thus the orbit visits \(Q(u1\,0^k,0,k)\) for every positive multiple \(k\) of three.

\begin{lemma}\label{leftheavy}
For positive multiples \(k,p\) of three, \(Q(u1\,0^{k+p},0,k)\) halts.
\end{lemma}
\begin{proof}
After \(k/3\) up sweeps and the zero exit the machine is at \(T(u1\,0^p,k)\). If \(p=3\), the entry reaches \(u10001\Bleft\) on the left. The four instructions \(\state{B1,E0,B0,C0}\) put D on the first one after \(u\).

If \(p\ge6\), put \(h=u1\,0^{p-4}\). The four-zero entry and \(k+1\) down sweeps give \(Q(h1\,0^{k+2},0,k+2)\). Another \(k/3\) up sweeps leave \(Q(h100,k,2)\). More generally \(Q(h100,b,2)\) halts for every \(b\). Put \(j=h100(10)^{b+1}\). A step \(\state{D0}\) and \(b+1\) copies of \(\state A01\run01\state A\) give \(j1\state A1\). The right exterior is blank, so \(\state{A1,C0,D0}\) give \(T(j1,0)\). Entry, exit and the next common entry position give
\[
 T(j1,0)\run P(j0,0,1)=Q(j0,0,1)\run T(j0,1)
                    \run j01\Bleft(10)01 .
\]
Move left across \(01\) by \(\state{B1,E0}\), then across \(b+1\) copies of \(10\) by \(\state{B0,C1}\). Now \(h100\Bleft\) is exposed. The steps \(\state{B0,C0}\) leave D on the displayed one.
\end{proof}

Cut the orbit words just before D. The longer prefix joined to the shorter suffix is the word of Lemma~\ref{leftheavy}, proving ordinary inseparability for both starts.

\subsection{The A start and the 176-state separator}\label{positive}

From A, the blank orbit reaches \(R(2)\) at time 14, where \(R(m)=T(\eps,m)\). For \(m\equiv2\pmod3\), entry, down sweeps, up sweeps and exit give
\[
 R(m)\run Q(1\,0^{m+2},0,m+2)\run Q(10,m+1,1)\run R(m+3).
\]
These diagonal lengths are \(1\) modulo three. We also use the gap repair \cite[Lemma~3]{prior}. For \(g,m\equiv1\pmod3\) and \(g\ge4\), two successive entry, sweep and exit cycles give
\begin{equation}\label{repair}
 T(1\,0^g,m)\run T(1\,0^{g-3},m+3).
\end{equation}
At \(g=1\) the word \(T(10,m)\) is already \(R(m+1)\).

Write \(D\) for the canonical language of the published 176-state partial DFA \cite{dfa}, and \(V\) for Blanchard's weighted invariant of this machine \cite{weights}. Call a configuration \emph{safe} if its run never reaches an undefined instruction. Define
\begin{equation}\label{gap}
 G_{ij}=1\,0^{4+3i}\,1(01)^{3j}0\,\state B1\,01,\qquad i,j\ge0 .
\end{equation}

\begin{lemma}\label{exit}
The DFA accepts blank and rejects every canonical word with head \(\state D1\). The configurations it accepts whose defined successors it rejects are exactly the \(G_{ij}\). Every \(G_{ij}\) is safe and lies outside \(V\).
\end{lemma}

Blank lies in \(D\cap V\). The weighted certificate is closed under defined steps, and Lemma~\ref{exit} puts every source of an exit from \(D\) outside \(V\). Thus \(D\cap V\) is closed and \(D\) contains the blank orbit. A halting path from \(D\) would have a first exit because its terminal word is rejected. Its exit source is safe, a contradiction. This proves ordinary separation. One rejecting sink completes the DFA.

\begin{proof}[Proof of Lemma~\ref{exit}]
Give each padded word the DFA value of its canonical form, and use a second automaton to recognize the padded forms of \eqref{gap}. Check blank acceptance and rejection at the undefined instruction directly. For each of the nine defined instructions and each adjacent bit \(y\), run copies of the padded DFA on the source and successor of the corresponding rewrite, and run the family recognizer on the source.
\[
\begin{array}{ll}
 l\,q x\,y r\ \longmapsto\ l b\,q' y\,r &\text{for a right move writing \(b\) into \(x\)},\\
 l y\,q x\,r\ \longmapsto\ l\,q' y\,b r &\text{for a left move writing \(b\) into \(x\)} .
\end{array}
\]
The product reads the common binary contexts \(l,r\) synchronously. Its reachable states exclude an accepted source outside the family with a rejected successor. Padding supplies the adjacent blank at either margin.

For the reverse inclusion, in the published DFA the prefix \(10000\) reaches state 39, where \(000\) loops. Reading the next \(1\) gives state 62, and \(010101\) sends each of states 62 and 174 to 174. From either state the suffix \(0\state B101\) ends in accepting state 12. The \(\state{B1}\) instruction replaces that suffix by \(\state E0001\), which ends in the rejecting sink. This checks every \(i,j\), including \(j=0\).

To prove safety, start from \(T(1\,0^{4+3i},3j+1)\). Applying \eqref{repair} \(i+1\) times reaches \(T(10,3(i+j)+4)=R(3(i+j)+5)\), on the nonhalting blank orbit. The first three steps from the same starting configuration, \(\state{A0,B0,C1}\), reach \(G_{ij}\). Determinism therefore makes \(G_{ij}\) safe.

On \(G_{ij}\) the left reader starts in state 0, reads the initial one into state 1 and the \(4+3i\) zeros into state 11. The next one reaches state 18, where \((01)^3\) loops. The last zero before the head gives state 26. The right reader, reading the reversal of \(01\), reaches state 2. The guard for \((\state B,1,26,2)\) is absent, so every \(G_{ij}\) lies outside \(V\), independently of the accumulated weights.
\end{proof}

For completeness, the obstruction to \emph{inductive} separation uses two pumps \cite[Theorem~1]{prior}. For any given DFA choose \(N\) and a multiple \(p\ge6\) of three such that reading \(0^t\) or \(0^{t+p}\) induces the same state transformation for every \(t\ge N\). Choose an orbit diagonal \(Q(1\,0^k,0,k)\) with \(k\ge N+4\), \(k\equiv1\pmod3\), and put \(s=k+p-4\). Write \(\sim\) for the two words reaching the same DFA state. Then
\begin{align*}
 Q(1\,0^k,0,k)&\sim Q(1\,0^{k+p},0,k)\run Q(10001\,0^s,0,s),\\
 Q(10001\,0^s,0,s)&\sim Q(10001\,0^{s-p},0,s)\run\text{a halt at }\state D1 .
\end{align*}
For the first run, \((k-1)/3\) up sweeps leave \(Q(1\,0^{p+1},k-1,1)\), and its exit gives \(T(1\,0^{p+1},k)\). Apply \eqref{repair} \((p-6)/3\) times to reach gap seven, then perform its four-zero entry and down sweeps. This exposes the second diagonal with \(s\) ones. In the second run, \((s-p)/3\) up sweeps give \(Q(10001,s-p,p)\). The next common up position exposes \(10001\Bleft\), which halts by the four left moves in Lemma~\ref{leftheavy}. Closure would force acceptance to propagate from the first pump to the second site and then to a halt. This proves the remaining assertion of Theorem~\ref{main}.

\section{The cost of separating bounded inputs}\label{width}

Call an automaton correct through width \(W\) if it accepts every word of \(\orb\) of span at most \(W\) and rejects every word of \(\halts\) of span at most \(W\). Only the input configuration is restricted. A computation from such an input may visit arbitrarily wider configurations.

\begin{theorem}\label{widththm}
For each of the sixteen ordinarily inseparable tables, an \(n\)-state DFA or \(\eps\)-NFA fails by the width in Table~\ref{widths}. In particular, every automaton correct through \(W\ge11\) has at least \(1+\lfloor(W-11)/8\rfloor\) states.
\end{theorem}
\begin{table}[ht]
\centering
\caption{Sufficient span for a violation by an \(n\)-state automaton.}\label{widths}
\begin{tabular}{@{}ll@{}}
\toprule
Family and blank start & Span in tape cells\\
\midrule
Counter-balanced A or D & \(4n+2\) or \(8n+4\)\\
Counter-inverting & \(\max(1,4n)\)\\
Cubic helix B or E & \(6n+9\) or \(6n+11\)\\
\bottomrule
\end{tabular}
\end{table}
\begin{proof}
Use indices \(0\le i\le n\) at the preceding cuts. Both fragment lengths are nondecreasing in \(i\). For the balanced machines the two banks grow, and for the inverters the lengths are \(3i\) and \(1+\max(1,i)\), including the scanned blank at \(i=0\). For the B- and E-start helices use \(k_i=3(i+1)\). Reflection sends a counter-balanced or helix diagonal \(L_i M R_i\) with fixed marked cell \(M=\state D0\) to \(R_i^{\rm rev}M'L_i^{\rm rev}\), with renamed mark \(M'\). Cutting after \(M'\) exchanges the banks and yields the reflected fatal splice on the literal alphabet.

By Lemma~\ref{cut} the last diagonal bounds every mixed span. In the balanced cases its span is \(2(e+(e+1)n)+2(e+1)n+2=4(e+1)n+2e+2\). The inverter span is \(3n+\max(1,n)=\max(1,4n)\). The helix span is \(|u|+3+2k_n=6n+9+|u|\), with \(|u|=0\) or \(2\). All thresholds are at most \(8n+11\), so correctness through \(W\) requires \(8n+11>W\).
\end{proof}

The DFA for \(\#4817065\) works at every span. For the other sixteen machines, Appendix~\ref{upper} gives a matching deterministic upper bound, hence the minimum is \(\Theta(W)\) for both DFAs and \(\eps\)-NFAs. The bounded-induction problem in \cite{cost} instead restricts the span of successor closure and requires global ordinary separation.

\section{An application to seven-state machines}\label{seven}

A seven-state restoring detour redirects one or more instructions of its five-state \emph{shadow} through two helper states. On return the tape, head and core state agree with the shadow step. If all calls return on finite-support tapes, both blank runs and halting runs of the shadow lift. Restricting a seven-state separator to core-state words and renaming the states would then separate the shadow. Toggles always return, and scans of ones eventually encounter a zero \cite{detours}. A scan of zeros needs a one in its direction of travel. For such a machine it suffices to show that every call on the particular orbit and halting witnesses returns.

At least 494 of the 721 distinct rows in the restoring-detour census of a historical 15,743,264-row list admit no ordinary regular separator (Table~\ref{liftcounts}). The 29 Finned \#3 and 20 Skelet \#17 lifts were established in \cite{detours}. Twelve rows admit both scan types and enter the count once through the scan of ones.

\begin{table}[ht]
\centering
\caption{Seven-state rows with an ordinary obstruction, by five-state shadow.}\label{liftcounts}
\begin{tabular}{@{}lrrr@{}}
\toprule
Shadow & Always-returning helper & Witness-specific zero scan & Total\\
\midrule
Finned \#3 & 27 & 2 & 29\\
Skelet \#17 & 20 & 0 & 20\\
B- and E-start cubic helices & 40 & 1 & 41\\
Counter-balanced & 262 & 1 & 263\\
Counter-inverting & 41 & 2 & 43\\
Four other Finned machines & 98 & 0 & 98\\
\midrule
Distinct rows & 488 & 6 & 494\\
\bottomrule
\end{tabular}
\end{table}

\paragraph{The cut for Skelet \#17.}
The obstruction due to dyuan01 pairs the prefix of one blank-reached checkpoint with a longer suffix from a later checkpoint \cite{dyuan}. To locate this cut on the actual tape, use savask's encoding. A tuple \((a_\ell,\ldots,a_0)\) spells
\[
   (10)^{a_\ell}\,1(10)^{a_{\ell-1}}\cdots 1(10)^{a_0}\,\state B0 .
\]
The B head scans the exterior blank immediately after the final two-bit block \cite[pp.~1--2]{savask}. Each extra copy of \(10\) changes the canonical word.

Take the blank-reached tuple just after the \emph{Zero} transition from \(E''\) in Xu's proof \cite[\S4, p.~9]{xu}. For \(k\ge1\), its two leading coordinates are zero, its Gray counter is \(n_k=2^{2k+1}-1\), its counting direction is \(-1\), and its terminal coordinate is \(z_k=n_k-4\). At \(E''\) the leading coordinate is 2 and the final one is \(z_k+1>0\), so this Zero transition is savask's physical rule (E1) \cite[Remark~1.1]{xu}. Split immediately before \((10)^{z_k}\state B0\).

Join the prefix for \(k\) to the suffix for \(j>k\). The terminal coordinate becomes \(z_j>n_k\). The Gray counter depends only on the other coordinates, and the parity of the full tuple determines the direction. Since both terminal coordinates are odd, the hybrid still counts down from \(n_k\). After \(t<n_k\) such decrements the counter is \(n_k-t>0\) and the terminal coordinate is \(z_j-t>0\). With negative direction Xu's rules force another \emph{Increment} step, decreasing both by one and preserving the direction. Induction gives exactly \(n_k\) decrements to counter zero. The two leading coordinates, at indices \(2k+2\) and \(2k+3\), receive zero increments because \(n_k/2^i<1/2\) for both indices \(i\) in Xu's nearest-integer formula \cite[Prop.~2.2]{xu}. The other Gray coordinates start positive, never decrease, and are even at counter zero. The terminal coordinate \(z_j-n_k\) is also positive and even, so savask's tape rule (E3) gives a halt \cite[p.~3]{savask}. Every diagonal is reached and every such hybrid halts, so Lemma~\ref{cut} supplies the obstruction inherited by the twenty rows.

\paragraph{The Finned shadows.}
For Finned \#3, Iijil proved ordinary inseparability by following the discrepancy between two block lengths and credited savask for the language argument \cite{iijil}. Blanchard identified all five Finned machines as provable CTL holdouts \cite{finnedsummary}. For the other four, the companion calculation \emph{Two counter phases for the five Finned machines} gives the literal returns. Each blank orbit visits all \(1^a0\,1^bq0\) with \(b=2a+c\), for a machine-specific constant \(c\). An odd discrepancy \(b-2a-c\) halts, and an even discrepancy becomes its negative half on the next return. Hence every nonzero discrepancy eventually halts. Lemma~\ref{cut} applies to each shadow and the transfer gives 71 toggle rows and 27 scans of ones.

\paragraph{The six witness-specific returns.}
For the two Finned \#3 zero scans, the witness calculations in \cite{detours} give a one to the right of each diverted instruction. Row 14,424,496 diverts only \(\state{E1}\) in the shadow. Row 15,538,865 diverts both \(\state{C0}\) and \(\state{E1}\) (literal seven-state \(\state{C0}\) and \(\state{G1}\)). At \(\state{C0}\) the immediately adjacent one was just written by \(\state{B0}\). At \(\state{E1}\) the extra one at the far end survives the sweep. Thus both callers return on the required paths.

For E-start helix row 3,192,869 set \(u=10\) in the second row of the blank return \eqref{secondreturn}. It visits \(Q(1\,0^{k+2},0,k+1)\) for positive multiples \(k\) of three. Keep the shorter left bank and take the longer right bank from the checkpoint indexed by \(l>k\). After \(k/3\) up sweeps, the hybrid reaches \(Q(100,k,l-k+1)\). At least four ones remain, so the next sweep puts B on the final zero of \(100\), and \(\state{B0,C0}\) reach \(\state{D1}\). The diverted \(\state{B1}\) moves left in base coordinates. Each such call in the splice lies in the right-scan tile, where \(\state{A0}\) has just written a one to its left. On the blank path the left scans retain a prefix ending in one, and the startup and special entry have a left delimiter as well. The full sweep-by-sweep check is in the companion seven-state calculation. This right-heavy splice supplies a halting witness on the literal seven-state table.

Lean checks the other three zero-scan rows directly. For 9,994,819 a surviving one bounds each required helper scan. Row 9,932,146 simulates the differing instruction in three steps. For 12,767,546 the cut can be made inside a helper excursion or at a core checkpoint.

\appendix
\section{The matching upper bound from the certificates}\label{upper}

For each of the seventeen weighted certificates and every \(W\ge1\), a deterministic finite automaton with \(O(W)\) states accepts the whole blank orbit and agrees with the certified invariant on every canonical word through span \(W\). The implied constant is specific to the certificate. Each certificate reads the left tape \(l\) forwards and the right tape \(r\) backwards, starting at state zero in each reader. The head state and bit and the two final reader states select a guard on the sum of the two weights \cite[\S4.7]{bb5}.

Fix the head state, scanned bit and final left-reader state. Let \(G\) consist of the right endpoints with a guard for this choice. Write \(z\) for the reversal of the right prefix read so far. A map entry \(M(p)=(q,d)\) records that the run from \(p\) on \(z\) ends at \(q\in G\). Its residual \(d\) is the run weight minus the least weight of any such run. The register \(R\) stores that minimum plus the left weight. Initially \(R\) is the left weight and \(M(p)=(p,0)\) exactly for \(p\in G\).

Appending a tape bit \(b\) changes \(z\) to \(bz\). For a right-reader edge \(p\xrightarrow{b/c}t\) of weight \(c\), form the temporary map \(\widetilde M\), then normalize its live entries by their least second component \(\mu\).
\[
 M(t)=(q,d)\ \Longrightarrow\ \widetilde M(p)=(q,c+d),
 \qquad
 \widetilde M(p)=(q,v)\ \Longrightarrow\ M'(p)=(q,v-\mu),\quad R'=R+\mu .
\]
An absent entry at \(t\) gives an absent entry at \(p\). An empty map stays empty and rejects. On every live run the update preserves the identity \(R+d=\text{left weight}+\text{right run weight}\). At the end, if \(M(0)=(q,d)\), test \(R+d\) against the guard for \(q\). An absent entry rejects.

Closing the initial maps under both input bits gives 8 distinct maps for each of the first ten certificate records, and \(15,17,15,36,606,881,1626\) for the remaining seven in file order. The companion program \word{check-weighted-suffixes.py} verifies these counts and that every live residual is at most one. Thus the maps and head data require only a fixed number of states per certificate.

If \(C\) is the largest absolute edge weight, then \(|R|\) is at most \(C\) times the number of unmarked cells read, both while reading the left tape and after adding the minimum right-run weight. Store the \(2CW+1\) values in \([-CW,CW]\). Overflow implies that the input exceeds span \(W\), so the automaton can accept on overflow subject only to the fixed syntax check. It accepts every certified word and agrees with the certificate throughout the promised span.

\section*{Verification and data}

Unpack \href{https://github.com/maxilism/busy-beaver-research/raw/refs/heads/main/support/verification/bb-5-2/regular-separation/lean-proof.zip}{\word{lean-proof.zip}} and run \word{lake build} with Lean 4.36.0 in its \word{transfers/formal\_review} directory. The principal theorems are listed below.
\begin{center}\small
\begin{tabular}{@{}ll@{}}
\toprule
Module & Theorems \\
\midrule
\word{SeventeenClassification} & \word{classification}, \word{publishedSeparator}\\
\word{PositiveSafetyConclusion} & \word{escape\_iff}\\
\word{WidthCases} & \word{dfaStateLower}, \word{enfaStateLower}\\
\bottomrule
\end{tabular}
\end{center}
The axiom report for the aggregate classification lists propositional extensionality, classical choice and quotient soundness.

The archive includes the original DFA and weighted certificate, the programs comparing their entries with the formal tables, and the catalog check on all seventeen distinct machines \cite{coq,weights,dfa}. Separate finite computations reproduce the suffix-map closure in Appendix~\ref{upper} and check the seven-state census against the published manifests. The accompanying seven-state and Finned calculations supply the transfer premises in Section~\ref{seven}, and Lean checks the three row-specific obstructions identified there. The Python replay \word{check\_skelet17.py} runs the literal five-state table. The \(k=1,2\) Skelet checkpoints occur after 5,558 and 576,105 steps, and the respective splices using the next suffix halt after 985 and 18,219 defined steps. The archive README gives the commands.

\begin{thebibliography}{99}\small\raggedright
\setlength{\itemsep}{1pt}\setlength{\parskip}{0pt}
\bibitem{bb5} The bbchallenge Collaboration, \emph{Determination of the fifth Busy Beaver value}, Proceedings of STOC 2026, 565--574. \href{https://doi.org/10.1145/3798129.3806376}{Conference paper}. \href{https://arxiv.org/abs/2509.12337v2}{Expanded version}, \S\S4.2, 4.6.2, 4.7.
\bibitem{coq} The bbchallenge Collaboration, \emph{Coq-BB5}, hardcoded parameters for the weighted finite automata verifier, commit \word{9142e21}. \href{https://github.com/ccz181078/Coq-BB5/blob/9142e219229baf2245d3f70851947230ea28a318/CoqBB5/BB5/BB5_Deciders_Hardcoded_Parameters/Verifier_WFAR_Hardcoded_Certificates.v}{Certificate file}.
\bibitem{weights} Iijil and contributors, \emph{MITMWFAR}, \word{SolvedFullCerts.txt}, commit \word{3937731}. \href{https://github.com/Iijil1/MITMWFAR/blob/3937731f9e1d801d6f6d77cf44787f93fc1ffc76/SolvedFullCerts.txt}{The seventeen tables and certificates}. \href{https://github.com/Iijil1/MITMWFAR/pull/1}{Blanchard's helix certificates}.
\bibitem{georgiev} G. Georgiev, \emph{Busy Beaver nonregular machines for class TM(5)}, 2003. \href{https://skelet.ludost.net/bb/nreg.html}{Machine catalog}.
\bibitem{inverter} G. Georgiev, \emph{Machine 5502432}, shift rules and proof of nonhalting. \href{https://skelet.ludost.net/bb/proof_5502432.txt}{Original calculation}.
\bibitem{holdouts} J. Blanchard, \emph{The 30 to 34 CTL holdouts from BB(5)}, April 9, 2023. \href{https://discuss.bbchallenge.org/t/the-30-to-34-ctl-holdouts-from-bb-5/141/1}{bbchallenge discussion and catalog}.
\bibitem{deka} K. Deka, observations on \(\#8210683\), reproduced by C. Xu, \emph{Discussion of machine 8210683}, 2023. \href{https://discuss.bbchallenge.org/t/8210683-does-hot-halt/163/1}{bbchallenge discussion}.
\bibitem{birget} J.-C. Birget, \emph{Intersection and union of regular languages and state complexity}, Information Processing Letters 43 (1992), 185--190. \href{https://doi.org/10.1016/0020-0190(92)90198-5}{Article}.
\bibitem{dyuan} dyuan01, irregularity argument for Skelet \#17, August 2, 2024. \href{https://discord.com/channels/960643023006490684/1082035052961091705/1269022847502647357}{Original discussion}. \href{https://wiki.bbchallenge.org/wiki/Skelet_17#Irregularity}{Busy Beaver Wiki account}.
\bibitem{savask} savask, \emph{Analysis of Skelet's machine 17}, March 5, 2023. \href{https://docs.bbchallenge.org/other/skelet17_savasks_analysis.pdf}{Tape encoding and transition rules}.
\bibitem{xu} C. Xu, \emph{Skelet \#17 and the fifth Busy Beaver number}, 2024. \href{https://arxiv.org/abs/2407.02426v2}{arXiv:2407.02426v2}.
\bibitem{iijil} Iijil, \emph{10756090 (Finned \#3) is irregular}, March 17, 2023. \href{https://discuss.bbchallenge.org/t/10756090-finned-3-is-irregular/137}{Checkpoint proof}.
\bibitem{finnedsummary} J. Blanchard, \emph{BB5's \textquotedblleft finned\textquotedblright{} machines (summary)}, May 19, 2024. \href{https://discuss.bbchallenge.org/t/bb5s-finned-machines-summary/234}{Family catalog and proofs}.
\bibitem{prior} maxilism, \emph{A five-state Turing machine with a regular separator but no regular inductive separator}, November 24, 2026. \href{https://github.com/maxilism/busy-beaver-research/blob/ee693aa3127ca2aaa09a37f22e8b846fa6048914/results/bb-5-2/cubic-helix.pdf}{Manuscript}.
\bibitem{dfa} maxilism, \emph{Canonical separator for machine 4817065}, 2026. \href{https://github.com/maxilism/busy-beaver-research/blob/ee693aa3127ca2aaa09a37f22e8b846fa6048914/support/verification/bb-5-2/cubic-helix/dfa.json}{Transition table}. \href{https://github.com/maxilism/busy-beaver-research/tree/ee693aa3127ca2aaa09a37f22e8b846fa6048914/support/verification/bb-5-2/cubic-helix}{Construction and verification files}.
\bibitem{cost} maxilism, \emph{The state cost of bounded induction for a five-state cubic helix}, November 24, 2026. \href{https://github.com/maxilism/busy-beaver-research/blob/ee693aa3127ca2aaa09a37f22e8b846fa6048914/results/bb-5-2/cubic-helix-state-complexity.pdf}{Manuscript}.
\bibitem{detours} maxilism, \emph{Two tape-restoring detours in 721 seven-state machines}, November 24, 2026. \href{https://github.com/maxilism/busy-beaver-research/blob/682350182398fa879b9db9fe73a043fe10596b47/results/bb-7-2/restoring-detours.pdf}{Manuscript}. \href{https://github.com/maxilism/busy-beaver-research/tree/682350182398fa879b9db9fe73a043fe10596b47/support/data/bb-7-2/restoring-detours}{Census manifests}.
\end{thebibliography}
\smallskip\noindent{\small\copyright\ 2026 maxilism. Licensed under \href{https://creativecommons.org/licenses/by/4.0/}{Creative Commons Attribution 4.0 International}.}
\end{document}
